\documentclass[10pt,a4paper]{amsart}
\usepackage{amsmath,amssymb,amsthm,graphicx}
\usepackage{verbatim}
\newenvironment{reference}{\comment}{\endcomment}
%\newenvironment{reference}{}{}
\newenvironment{slogan}{\comment}{\endcomment}
\newenvironment{history}{\comment}{\endcomment}

% For commutative diagrams we use Xy-pic
\usepackage[all]{xy}

% We use 2cell for 2-commutative diagrams.
\xyoption{2cell}
% Theorem environments.
%
\theoremstyle{plain}
\newtheorem{theorem}[subsection]{Theorem}
\newtheorem{proposition}[subsection]{Proposition}
\newtheorem{lemma}[subsection]{Lemma}

\theoremstyle{definition}
\newtheorem{definition}[subsection]{Definition}
\newtheorem{example}[subsection]{Example}
\newtheorem{exercise}[subsection]{Exercise}
\newtheorem{situation}[subsection]{Situation}

\theoremstyle{remark}
\newtheorem{remark}[subsection]{Remark}
\newtheorem{remarks}[subsection]{Remarks}

\numberwithin{equation}{subsection}

% Macros
%
\def\lim{\mathop{\mathrm{lim}}\nolimits}
\def\colim{\mathop{\mathrm{colim}}\nolimits}
\def\Spec{\mathop{\mathrm{Spec}}}
\def\Hom{\mathop{\mathrm{Hom}}\nolimits}
\def\Ext{\mathop{\mathrm{Ext}}\nolimits}
\def\SheafHom{\mathop{\mathcal{H}\!\mathit{om}}\nolimits}
\def\SheafExt{\mathop{\mathcal{E}\!\mathit{xt}}\nolimits}
\def\Sch{\mathit{Sch}}
\def\Mor{\mathop{\mathrm{Mor}}\nolimits}
\def\Ob{\mathop{\mathrm{Ob}}\nolimits}
\def\Sh{\mathop{\mathit{Sh}}\nolimits}
\def\NL{\mathop{N\!L}\nolimits}
\def\CH{\mathop{\mathrm{CH}}\nolimits}
\def\proetale{{pro\text{-}\acute{e}tale}}
\def\etale{{\acute{e}tale}}
\def\QCoh{\mathit{QCoh}}
\def\Ker{\mathop{\mathrm{Ker}}}
\def\Im{\mathop{\mathrm{Im}}}
\def\Coker{\mathop{\mathrm{Coker}}}
\def\Coim{\mathop{\mathrm{Coim}}}

% Boxtimes
%
\DeclareMathSymbol{\boxtimes}{\mathbin}{AMSa}{"02}

%
% Macros for moduli stacks/spaces
%
\def\QCohstack{\mathcal{QC}\!\mathit{oh}}
\def\Cohstack{\mathcal{C}\!\mathit{oh}}
\def\Spacesstack{\mathcal{S}\!\mathit{paces}}
\def\Quotfunctor{\mathrm{Quot}}
\def\Hilbfunctor{\mathrm{Hilb}}
\def\Curvesstack{\mathcal{C}\!\mathit{urves}}
\def\Polarizedstack{\mathcal{P}\!\mathit{olarized}}
\def\Complexesstack{\mathcal{C}\!\mathit{omplexes}}
% \Pic is the operator that assigns to X its picard group, usage \Pic(X)
% \Picardstack_{X/B} denotes the Picard stack of X over B
% \Picardfunctor_{X/B} denotes the Picard functor of X over B
\def\Pic{\mathop{\mathrm{Pic}}\nolimits}
\def\Picardstack{\mathcal{P}\!\mathit{ic}}
\def\Picardfunctor{\mathrm{Pic}}
\def\Deformationcategory{\mathcal{D}\!\mathit{ef}}
\usepackage{cleveref}
\crefname{originalreference}{original reference}{original references}
\Crefname{originalreference}{Original reference}{Original references}
\makeatletter
\newlabel{lemma-algebraically-closed}{{09GR}{source}{}{sourcepointer-1}{}}
\newlabel{lemma-algebraically-closed@cref}{{[originalreference][0][]09GR}{[1][0][]source}{}{}{}}
\makeatother
\crefname{originalreference}{original reference}{original references}
\Crefname{originalreference}{Original reference}{Original references}
\makeatletter
\newlabel{lemma-characterize-normal}{{09HQ}{source}{}{sourcepointer-2}{}}
\newlabel{lemma-characterize-normal@cref}{{[originalreference][0][]09HQ}{[1][0][]source}{}{}{}}
\makeatother
\crefname{originalreference}{original reference}{original references}
\Crefname{originalreference}{Original reference}{Original references}
\makeatletter
\newlabel{lemma-map-into-algebraic-closure}{{09GU}{source}{}{sourcepointer-3}{}}
\newlabel{lemma-map-into-algebraic-closure@cref}{{[originalreference][0][]09GU}{[1][0][]source}{}{}{}}
\makeatother
\title{Existence, normality and uniqueness of the splitting field of a polynomial}
\author{The Stacks Project Authors}
\date{stacks-project:fields.tex:local-snapshot-20261012; GNU FDL 1.2 or later}
\begin{document}
\maketitle
\noindent\textit{Editorial note.} All delivered mathematical text is the author's own, copied from the pinned Stacks Project chapter file fields.tex; no statement or proof step was rewritten, reordered or removed.\par\medskip
\noindent\textit{Editorial note.} The counted proof is the author's complete argument: The uniqueness isomorphism sigma\textasciicircum{}\{-1\} composed with the inclusion realizes 'unique up to (nonunique) isomorphism' of the statement.\par\medskip
\noindent\textit{Editorial note.} The label lemma-splitting-field names the counted statement, and the proof environment immediately following it in the source is delivered verbatim as its proof.\par\medskip
\noindent\textit{Editorial note.} The counted statement is carried by the author's own lemma environment in the chapter source.\par\medskip
\noindent\textit{Allowed standard prerequisites (editorial).}\begin{itemize}
\item Same chapter: lemma-algebraically-closed (factorization over an algebraic closure), lemma-map-into-algebraic-closure (extension of embeddings), lemma-characterize-normal (normality via equal images of embeddings)
\end{itemize}\medskip
\ifdefined\hypertarget\hypertarget{sourcepointer-1}{}\fi
\noindent\textit{Source reference.} Original reference 09GR: sources/stacks-fields/pinned-original-source.tex, source lines 1128--1128. Prerequisite pointer invoked at 2276; it names a standard fact of the same chapter that the counted argument uses, and every step of the counted proof built on it is retained verbatim in the delivered body.\par\medskip
\ifdefined\hypertarget\hypertarget{sourcepointer-2}{}\fi
\noindent\textit{Source reference.} Original reference 09HQ: sources/stacks-fields/pinned-original-source.tex, source lines 2096--2096. Prerequisite pointer invoked at 2291; it names a standard fact of the same chapter that the counted argument uses, and every step of the counted proof built on it is retained verbatim in the delivered body.\par\medskip
\ifdefined\hypertarget\hypertarget{sourcepointer-3}{}\fi
\noindent\textit{Source reference.} Original reference 09GU: sources/stacks-fields/pinned-original-source.tex, source lines 1240--1240. Prerequisite pointer invoked at 2287; it names a standard fact of the same chapter that the counted argument uses, and every step of the counted proof built on it is retained verbatim in the delivered body.\par\medskip
\begin{lemma}
\label{lemma-splitting-field}
Let $F$ be a field. Let $P \in F[x]$ be a nonconstant polynomial.
There exists a smallest field extension $E/F$ such that $P$
splits completely over $E$. Moreover, the field extension $E/F$ is normal
and unique up to (nonunique) isomorphism.
\end{lemma}

\begin{proof}
Choose an algebraic closure $\overline{F}$. Then we can write
$P = c (x - \beta_1) \ldots (x - \beta_n)$ in $\overline{F}[x]$, see
Lemma \ref{lemma-algebraically-closed}. Note that $c \in F^*$. Set
$E = F(\beta_1, \ldots, \beta_n)$. Then it is clear that $E$ is
minimal with the requirement that $P$ splits completely over $E$.

\medskip\noindent
Next, let $E'$ be another minimal field extension of $F$ such that
$P$ splits completely over $E'$. Write
$P = c (x - \alpha_1) \ldots (x - \alpha_n)$ with $c \in F$ and
$\alpha_i \in E'$. Again it follows from minimality that
$E' = F(\alpha_1, \ldots, \alpha_n)$. Moreover, if we pick
any $\sigma : E' \to \overline{F}$
(Lemma \ref{lemma-map-into-algebraic-closure})
then we immediately see that $\sigma(\alpha_i) = \beta_{\tau(i)}$
for some permutation $\tau : \{1, \ldots, n\} \to \{1, \ldots, n\}$.
Thus $\sigma(E') = E$. This implies that $E'$ is a normal extension
of $F$ by Lemma \ref{lemma-characterize-normal}
and that $E \cong E'$ as extensions of $F$ thereby finishing the proof.
\end{proof}

\end{document}
